% Complete analytical appendix proposal. Requires amsmath, amssymb, natbib.
% Numerical claims and a final figure caption must follow the audit note.
% Scope: linear prediction and finite fixed random-feature networks only.
\section{Target-dependent fitting: derivations and nonlinear features}
\label{app:target_fitting}

\subsection{Data, targets, and optimization time}
Draw $n$ clean examples independently from a population $P$ and keep this
dataset fixed. At a fixed flow time $t\in(0,1)$, form
\[
 \mathbf z=t\mathbf x+(1-t)\boldsymbol\epsilon,
 \qquad \boldsymbol\epsilon\sim\mathcal N(0,I_d),
 \qquad \boldsymbol\epsilon\perp\mathbf x.
\]
The training expectation selects a stored clean example uniformly and
averages over fresh corruption. The population expectation draws both a
fresh clean example and fresh corruption. Write these expectations as
$\mathbb E_{\mathrm{tr}}$ and $\mathbb E_P$.
The native targets are
$\mathbf y_x=\mathbf x$, $\mathbf y_\epsilon=\boldsymbol\epsilon$, and
$\mathbf y_v=\mathbf x-\boldsymbol\epsilon$.

Let $\boldsymbol\phi:\mathbb R^d\to\mathbb R^p$ be fixed and train only
$B_q\in\mathbb R^{d\times p}$ in
$f_q(\mathbf z)=B_q\boldsymbol\phi(\mathbf z)$. Define
\begin{align}
 U&=\mathbb E_{\mathrm{tr}}[\boldsymbol\phi\boldsymbol\phi^\top],&
 V_q&=\mathbb E_{\mathrm{tr}}[\boldsymbol\phi\mathbf y_q^\top],\\
 C&=\mathbb E_P[\boldsymbol\phi\boldsymbol\phi^\top],&
 V_{P,q}&=\mathbb E_P[\boldsymbol\phi\mathbf y_q^\top].
\end{align}
Assume these moments and the displayed risks are finite.
The feature moments $U,C$ have shape $p\times p$; the feature--target
moments $V_q,V_{P,q}$ have shape $p\times d$.
With $c_{q,r}=d^{-1}\mathbb E_r\|\mathbf y_q\|^2$, the losses are
\begin{align}
 \mathcal L_{q,\mathrm{tr}}(B)
 &=c_{q,\mathrm{tr}}+\frac1d\operatorname{tr}(BUB^\top)
   -\frac2d\operatorname{tr}(BV_q),\\
 \mathcal L_{q,P}(B)
 &=c_{q,P}+\frac1d\operatorname{tr}(BCB^\top)
   -\frac2d\operatorname{tr}(BV_{P,q}).
\end{align}
Consequently $\nabla_B\mathcal L_{q,\mathrm{tr}}
=(2/d)(BU-V_q^\top)$. Throughout this appendix use the rescaled
gradient-flow time for which
\begin{equation}
 \dot B_q=V_q^\top-B_qU.
 \label{eq:app_readout_flow}
\end{equation}
For gradient descent with step size $\eta$ on the stated per-coordinate
loss, this clock is $\tau=2k\eta/d$. Discrete or stochastic updates
approximate this flow; the matrix solution below is exact for the flow.

\subsection{Exact solution for any fixed feature map}
For $U\mathbf u_j=\lambda_j\mathbf u_j$, solve
$\frac{d}{d\tau}(B_q\mathbf u_j)
=V_q^\top\mathbf u_j-\lambda_jB_q\mathbf u_j$. This gives
\begin{equation}
 B_q(\tau)=B_q(0)e^{-\tau U}+V_q^\top g_\tau(U),\qquad
 g_\tau(\lambda)=\frac{1-e^{-\lambda\tau}}{\lambda},
 \label{eq:app_readout_solution}
\end{equation}
where $g_\tau(0)=\tau$. The zero-eigenvalue convention makes the formula
valid for singular $U$. In a training-null direction,
$V_q^\top\mathbf u_j=0$ because
$\mathbb E_{\mathrm{tr}}(\mathbf u_j^\top\boldsymbol\phi)^2=0$.
The corresponding initial component remains unchanged.

An eigenvalue $\lambda_j>0$ gives a response time $1/\lambda_j$.
These response times are shared across targets. The projection
$V_q^\top\mathbf u_j$ determines the target-dependent response amplitude.
Population error also uses $C$ and $V_{P,q}$; eigenvalues of $U$ alone
do not determine the error curves.

\subsection{Velocity errors and the change in the gap}
For any feature map, convert predictions to velocities using
\[
 \widehat{\mathbf v}_x=\frac{B_x\boldsymbol\phi-\mathbf z}{1-t},\qquad
 \widehat{\mathbf v}_\epsilon=\frac{\mathbf z-B_\epsilon\boldsymbol\phi}{t},
 \qquad \widehat{\mathbf v}_v=B_v\boldsymbol\phi.
\]
These conversions imply the exact samplewise error identities
$\widehat{\mathbf v}_x-\mathbf v=(f_x-\mathbf x)/(1-t)$ and
$\widehat{\mathbf v}_\epsilon-\mathbf v
=-(f_\epsilon-\boldsymbol\epsilon)/t$.
Therefore
\begin{gather*}
 R^v_{q,r}=\frac1d\mathbb E_r\|\widehat{\mathbf v}_q-\mathbf v\|^2
 =\delta_q^2\mathcal L_{q,r},\\
 (\delta_x,\delta_\epsilon,\delta_v)=((1-t)^{-1},t^{-1},1).
\end{gather*}
Define $\mathcal G_q=R^v_{q,P}-R^v_{q,\mathrm{tr}}$ and
$\Delta\mathcal G_q(\tau)=\mathcal G_q(\tau)-\mathcal G_q(0)$.
A positive change indicates an increasing training advantage.
Under exact gradient flow, training error cannot increase.
Population error can decrease or increase; the gap alone does not determine
that change.
For $B_q(0)=0$,
\begin{equation}
 \Delta\mathcal G_q
 =\frac{\delta_q^2}{d}\left\{
 \operatorname{tr}[B_q(C-U)B_q^\top]
 -2\operatorname{tr}[B_q(V_{P,q}-V_q)]\right\}.
 \label{eq:app_native_gap}
\end{equation}
The target second-moment constants cancel on subtracting the initial gap.

\subsection{Linear predictors share one velocity equation}
Set $\boldsymbol\phi(\mathbf z)=\mathbf z$, so $p=d$, and write
$S=n^{-1}\sum_i\mathbf x_i\mathbf x_i^\top$ and
$\Sigma=\mathbb E_P[\mathbf x\mathbf x^\top]$.
These are second moments, without subtracting the mean. Noise averaging gives
\begin{align}
 U&=t^2S+(1-t)^2I_d,& C&=t^2\Sigma+(1-t)^2I_d,\\
 V_x&=tS,&V_{P,x}&=t\Sigma,\\
 V_\epsilon&=(1-t)I_d,&V_{P,\epsilon}&=(1-t)I_d,\\
 V_v&=tS-(1-t)I_d,&V_{P,v}&=t\Sigma-(1-t)I_d.
\end{align}
Since $0<t<1$, $U$ is positive definite even if $S$ is singular.
The linear conversion matrices are those in
Eq.~\eqref{eq:main_velocity_conversion}.
The identities $U=tV_x+(1-t)V_\epsilon$ and $V_v=V_x-V_\epsilon$
give, for every target,
\begin{equation}
 \dot A_q=V_v^\top-A_qU,
 \qquad
 A_q(\tau)=A_q(0)e^{-\tau U}+V_v^\top g_\tau(U).
 \label{eq:app_velocity_solution}
\end{equation}
All three converge to $A_\star=V_v^\top U^{-1}$.
Matching an arbitrary initial velocity matrix $A_0$ requires
\begin{equation}
 \begin{aligned}
 B_x(0)&=I_d+(1-t)A_0,\\
 B_\epsilon(0)&=I_d-tA_0,\qquad B_v(0)=A_0.
 \end{aligned}
 \label{eq:app_matched_initialization}
\end{equation}
Their velocity trajectories then coincide. The same statement holds for
discrete linear updates with the same step size and the same corruption
draws: the conversion identities hold for each batch.

In velocity coordinates, with $A_{q,0}=A_q(0)$, the gap change is
\begin{align}
 \Delta\mathcal G_q
 ={}&\frac1d\operatorname{tr}\left[
 A_q(C-U)A_q^\top-A_{q,0}(C-U)A_{q,0}^\top\right]\nonumber\\
 &-\frac2d\operatorname{tr}[(A_q-A_{q,0})(V_{P,v}-V_v)].
 \label{eq:app_velocity_gap}
\end{align}
Identical limiting velocity maps imply identical limiting risks and raw
gaps. The gap changes can retain different limits for one realized dataset,
because the subtracted initial gaps can differ.

\subsection{Early-time coefficients and their assumptions}
First take any population with finite fourth moment and zero native
initialization. Define
\begin{equation}
 \kappa=\frac{\mathbb E\|\mathbf x\|^4
                  -\operatorname{tr}(\Sigma^2)}{nd}
 =\frac1d\mathbb E_{\mathcal D}\operatorname{tr}[(S-\Sigma)^2].
 \label{eq:app_kappa}
\end{equation}
Independence of the clean draws gives
\[
 \mathbb E\operatorname{tr}(S^2-S\Sigma)=d\kappa,
 \qquad \mathbb E(S-\Sigma)=0.
\]
Substituting $B_q(\tau)=\tau V_q^\top+O(\tau^2)$ into
Eq.~\eqref{eq:app_native_gap} yields the initial expected slopes
\begin{equation}
 \left.\frac{d}{d\tau}\mathbb E\Delta\mathcal G_q\right|_{\tau=0}
 =\begin{cases}
 2t^2\kappa/(1-t)^2,&q=x,\\
 0,&q=\epsilon,\\
 2t^2\kappa,&q=v.
 \end{cases}
 \label{eq:app_linear_slopes}
\end{equation}
The expectation is over datasets. The first derivative for noise prediction
vanishes for every dataset, since $V_\epsilon=V_{P,\epsilon}$.

For the higher-order calculation, assume $\mathbb E\|\mathbf x\|^6<\infty$.
Gaussian populations satisfy this sufficient condition for the stated
expected Taylor remainder.
The noise readout has the expansion
\[
 B_\epsilon=(1-t)\left(\tau I_d-\tfrac12\tau^2U
                            +\tfrac16\tau^3U^2\right)+O(\tau^4).
\]
As $C-U=t^2(\Sigma-S)$, Eq.~\eqref{eq:app_native_gap} becomes
\begin{align}
 \Delta\mathcal G_\epsilon(\tau)
 ={}&\frac{(1-t)^2}{d}\operatorname{tr}(\Sigma-S)\tau^2\nonumber\\
 &-\frac{(1-t)^2}{d}\operatorname{tr}[U(\Sigma-S)]\tau^3
 +O(\tau^4).
 \label{eq:app_noise_expansion}
\end{align}
The quadratic term vanishes in dataset expectation, not on each dataset.
Since $\mathbb E\operatorname{tr}[U(\Sigma-S)]=-t^2d\kappa$,
\begin{equation}
 \mathbb E\Delta\mathcal G_\epsilon(\tau)
 =t^2(1-t)^2\kappa\tau^3+O(\tau^4).
 \label{eq:app_noise_cubic}
\end{equation}
For $P=\mathcal N(0,I_d)$, $\kappa=(d+1)/n$.
Taking $t=1/2$ gives Eq.~\eqref{eq:main_early_gap}.
These leading coefficients scale as $1/n$ at fixed $d$; this does not
establish a copying-onset time proportional to $n$.

\paragraph{Population error at the Gaussian midpoint.}
For $P=\mathcal N(0,I_d)$ at $t=1/2$, $\mathbf z$ and $\mathbf v$ are
independent Gaussian vectors.
Thus $\mathbb E[\mathbf v\mid\mathbf z]=0$ and
\begin{equation}
 R^v_P(A)=2+\frac{\|A\|_F^2}{2d}.
 \label{eq:app_midpoint_population}
\end{equation}
The initial errors are $4,4,2$ for $x,\epsilon,v$.
The nonnegative excess $R^v_P-2$ directly measures deviation from the
population-optimal predictor. This example does not contain a population
learning phase for zero-initialized velocity prediction.

\paragraph{Effect of random native initialization.}
Suppose the entries of $B_q(0)$ are independent, have mean zero and variance
$\sigma^2$, and are independent of the dataset. Averaging over initialization
and datasets gives the initial slopes
\begin{align}
 a_x&=\frac{2t^2+2t^4d\sigma^2}{(1-t)^2}\kappa,\\
 a_\epsilon&=2t^2d\sigma^2\kappa,\\
 a_v&=(2t^2+2t^4d\sigma^2)\kappa.
 \label{eq:app_random_initialization_slopes}
\end{align}
To obtain the additional term, average the quadratic initial-readout
contribution to the derivative of the gap:
$-(2\delta_q^2/d)\mathbb E\operatorname{tr}
[B_q(0)U(C-U)B_q(0)^\top]$.
Use $\mathbb E[B_q(0)^\top B_q(0)]=d\sigma^2I_d$ and
$\mathbb E\operatorname{tr}[U(C-U)]=-t^4d\kappa$.
Noise prediction therefore has a linear expected gap term for nonzero
$\sigma$. A random-initialization comparison can test finite-time ordering;
it cannot establish robustness of the zero-initialization cubic law.

\subsection{Fixed nonlinear features: what extends and what changes}
For the nonlinear model, take
\begin{gather*}
 \boldsymbol\phi(\mathbf z)=\tanh(W\mathbf z/\sqrt d),
 \\ W\in\mathbb R^{p\times d},\qquad W_{ij}\sim\mathcal N(0,1),
\end{gather*}
and hold $W$ fixed independently of the dataset. Only $B_q$ is trained.
Equations~\eqref{eq:app_readout_solution} and
\eqref{eq:app_native_gap} remain exact, conditional on these features and
the dataset. Computing the moments by numerical integration introduces
integration error; it is separate from the exact matrix solution.

\paragraph{An exact expected initial slope.}
For each clean example define
\[
 H_q(\mathbf x)
 =\mathbb E_{\boldsymbol\epsilon}
 [\boldsymbol\phi(t\mathbf x+(1-t)\boldsymbol\epsilon)\mathbf y_q^\top
  \mid\mathbf x].
\]
Then $V_q=n^{-1}\sum_iH_q(\mathbf x_i)$ and
$V_{P,q}=\mathbb E_{\mathbf x}H_q(\mathbf x)$.
For zero native initialization, differentiation of
Eq.~\eqref{eq:app_native_gap} gives
\[
 \Delta\mathcal G_q'(0)
 =\frac{2\delta_q^2}{d}
   \bigl(\|V_q\|_F^2-\langle V_q,V_{P,q}\rangle_F\bigr).
\]
Because $\mathbb E_{\mathcal D}V_q=V_{P,q}$, averaging yields
\begin{equation}
 \mathbb E_{\mathcal D}\Delta\mathcal G_q'(0)
 =\frac{2\delta_q^2}{nd}
   \mathbb E_{\mathbf x}\|H_q(\mathbf x)-V_{P,q}\|_F^2.
 \label{eq:app_nonlinear_slope}
\end{equation}
This identity requires finite second moments of $H_q$ and holds
conditional on $W$. It shows that the initial expected gap depends on how
the feature--target relationship varies between clean examples.
It is nonnegative, but does not imply an ordering across targets.

Gaussian integration by parts gives
\begin{equation}
 H_\epsilon(\mathbf x)
 =(1-t)\mathbb E_{\boldsymbol\epsilon}
       [J_{\boldsymbol\phi}(\mathbf z)\mid\mathbf x],
 \label{eq:app_stein_features}
\end{equation}
where $J_{\boldsymbol\phi}$ is the $p\times d$ Jacobian.
For linear features this is the constant $(1-t)I_d$, explaining the zero
initial noise slope. For tanh features,
\[
 J_{\boldsymbol\phi}(\mathbf z)
 =\operatorname{diag}\!\left[1-\tanh^2(W\mathbf z/\sqrt d)\right]
   W/\sqrt d.
\]
Its conditional average can depend on $\mathbf x$.
Whenever this dependence has nonzero variance, noise prediction has a
positive linear expected gap term. The cubic linear-model result is
therefore not a general property of noise prediction.

\paragraph{The velocity function classes need not coincide.}
Let $\mathcal F$ be the linear space of functions $B\boldsymbol\phi$.
The decoded classes are
$\mathcal F-\mathbf z/(1-t)$,
$\mathcal F+\mathbf z/t$, and $\mathcal F$.
They coincide if $\mathbf z$ is representable by the features. A finite
tanh readout is bounded, so it cannot represent $\mathbf z$ exactly under
Gaussian corruption. The limiting velocity
predictors can therefore differ in addition to their initial functions.
Including the input as explicit linear features restores this property,
but defines a different model from the tanh model studied here.

For any fixed features, there is nevertheless one velocity function equation:
\begin{equation}
 \partial_\tau\widehat{\mathbf v}_q(\mathbf z)
 =\mathbb E_{\mathrm{tr}}\!\left[
 (\mathbf v'-\widehat{\mathbf v}_q(\mathbf z'))
 \boldsymbol\phi(\mathbf z')^\top\right]\boldsymbol\phi(\mathbf z).
 \label{eq:app_function_flow}
\end{equation}
Here $(\mathbf z',\mathbf v')$ is a fresh pair under the training expectation.
This follows directly by converting Eq.~\eqref{eq:app_readout_flow}.
If identical initial velocity functions are representable in all three
models, this common equation preserves their equality.

\paragraph{Computing population error.}
The eigenvalues of the feature matrix alone do not determine target
ordering. For example, with zero initialization, the population quadratic
term contains
\[
 \operatorname{tr}\!\left[V_q^\top g_\tau(U)C g_\tau(U)V_q\right],
\]
and the linear term contains
$\operatorname{tr}[V_q^\top g_\tau(U)V_{P,q}]$.
These expressions depend on the target moments and their relationships
with the empirical and population feature matrices. They can be evaluated
for a finite network once its moments are available. This gives an exact
readout solution conditional on those moments, without a limit in dimension
or width. Monte Carlo estimates of the moments and finite-step SGD add
separate numerical errors that must be measured when comparing the curves.
